% s02-1-basics.tex —— 2.1 节：D-mod(Bun_G) 的紧生成（底本印刷页 16–29）

\section{\texorpdfstring{$D\text{-}\mathrm{mod}(\Bun_G)$ 的紧生成}{D-mod(G) 的紧生成}}
\label{sec:cg-bung}

本节简要综述 Drinfeld 与 Gaitsgory 在 [DG15] 中已证的结论，即
$D\text{-}\mathrm{mod}(\Bun_G)$ 的紧生成。需要强调：本节的结果没有一项属于本文作者。

\srcpage{16}

\subsection{DG 范畴的紧生成}
\label{sec:cg-dg}

我们的设定是 $\DGCat_{\mathrm{cont}}$；也就是说，除非另有说明，DG 范畴都假定余完备，
范畴之间的函子都假定连续。
称对象 $c\in C$ 为\emph{紧的}，若函子
\[
\Maps_C(c,-)\colon C\to \Vect
\]
连续。

\begin{definition}[{[DG15, 1.2.2]}]
\label{def:comp-gen}
设 $C$ 为 DG 范畴。称 $C$ \emph{紧生成}，若存在一簇紧对象 $\{c_\alpha\in C\}_{\alpha\in I}$，
使得
\[
\Maps_C(c_\alpha,c')=0\quad\forall\alpha
\]
蕴涵 $c'=0$。
\end{definition}


下面两个命题是紧生成这一性质的稳定性。

\begin{proposition}
\label{prop:comp-gen-adjoint}
设 $C$、$D$ 为一对 DG 范畴，
\[
F\colon C\to D
\]
为它们之间的（连续）函子。则 $F$ 有（未必连续的）右伴随
\[
G\colon D\to C.
\]
若 $G$ 连续且保守，且 $C$ 紧生成，则 $D$ 也紧生成。
\end{proposition}

现考虑函子
\[
\Psi\colon I\to \DGCat_{\mathrm{cont}},
\]
它把
\[
(i\in I)\mapsto C_i,\qquad (i\to j\in I)\mapsto (\psi_{i,j}\colon C_i\to C_j).
\]
对每个 $i\to j\in I$，
\[
\psi_{i,j}\colon C_i\to C_j
\]
有右伴随
\[
\phi_{j,i}\colon C_j\to C_i.
\]
若 $\phi_{j,i}$ 对一切 $i\to j$ 都连续，则得到函子
\[
\Phi\colon I^{\mathrm{op}}\to \DGCat_{\mathrm{cont}},
\]
并且有事实（见 \cite[命题 1.7.5]{DG15}）：
\[
\colim_{i\in I}(C_i,\psi_{i,j})
\]
典范等价于
\[
\lim_{i\in I^{\mathrm{op}}}(C_i,\phi_{j,i}).
\]

\begin{proposition}[{[DG15, 推论 1.9.4]}]
\label{prop:comp-gen-colim}
在上述情形下，若每个 $C_i$ 都紧生成，则 $\colim_{i\in I}(C_i,\psi_{i,j})$ 紧生成。
\end{proposition}

\srcpage{17}

\subsection{QCA 叠}
\label{sec:qca}

本小节极其简略地综述 [DG13] 中与 $D$-模紧生成相关的部分。

先回顾 $\IndCoh$ 的定义。给定拟紧（局部 Noether、且几乎有限型）的 DG 概形 $Z$，可考虑
\[
\Coh(Z)\subset \QCoh(Z),
\]
并定义
\[
\IndCoh(Z):=\Ind(\Coh(Z)).
\]
给定映射 $f\colon Z_1\to Z_2$，有（连续）函子
\[
f^{!}\colon \IndCoh(Z_2)\to \IndCoh(Z_1),
\]
其定义见 [Gai]。这些数据提升为函子
\[
\IndCoh^{!}\colon (\DGSch_{\mathrm{aft}})^{\mathrm{op}}\to \DGCat_{\mathrm{cont}},
\]
它把
\[
(Z\in \DGSch_{\mathrm{aft}})\mapsto \IndCoh(Z),\qquad
(f\colon Z_1\to Z_2)\mapsto (f^{!}\colon \IndCoh(Z_2)\to \IndCoh(Z_1)).
\]
\srcpage{18}
记 $\mathrm{PreStk}$ 为（局部几乎有限型的）预叠范畴。可沿包含
\[
\DGSch^{\mathrm{op}}_{\mathrm{aft}}\to \mathrm{PreStk}^{\mathrm{op}}
\]
左 Kan 延拓 $\IndCoh^{!}$，得到仍记作
\[
\IndCoh^{!}\colon \mathrm{PreStk}^{\mathrm{op}}\to \DGCat_{\mathrm{cont}}
\]
的函子（记号略有滥用）。更具体地，
\[
\IndCoh(Y)=\lim_{Z\to Y}\IndCoh(Z),
\]
其中极限沿一切映入 $Y$ 的拟紧 DG 概形 $Z$ 取。换言之，预叠 $Y$ 上的一个归纳凝聚层，
等价于对每个从拟紧 DG 概形出发的映射 $g\colon Z\to Y$ 给出一族相容的
$g^{!}\mathcal{F}\in \IndCoh(Z)$。

现设 $Y$ 为代数叠。可考虑全子范畴
\[
\Coh(Y)\subset \IndCoh(Y),
\]
由满足「对一切从拟紧 DG 概形出发的光滑映射 $g\colon Z\to Y$ 都有
$g^{!}\mathcal{F}\in \Coh(Z)$」的 $\mathcal{F}\in \IndCoh(Y)$ 组成。包含
\[
\Coh(Y)\subset \IndCoh(Y)
\]
\srcpage{19}
给出函子
\[
\Ind(\Coh(Y))\to \IndCoh(Y);
\]
若 $Y$ 是拟紧 DG 概形，由 $\IndCoh$ 的定义，该函子是等价。但一般而言它不必是等价。

\begin{definition}[{[DG15, 定义 2.2.2]}]
\label{def:qca}
设 $Y$ 为代数叠。若 $Y$ 的几何点的自同构群都是仿射的，则称 $Y$ \emph{局部 QCA}。
若 $Y$ 拟紧且局部 QCA，则称 $Y$ \emph{QCA}。
\end{definition}

\begin{proposition}[{[DG13, 定理 3.3.5]}]
\label{prop:indcoh-coh}
设 $Y$ 为 QCA 代数叠。则上述函子
\[
\Ind(\Coh(Y))\to \IndCoh(Y)
\]
是等价。特别地，对 QCA 代数叠 $Y$，$\IndCoh(Y)$ 由 $\Coh(Y)$ 紧生成。
\end{proposition}

下面讨论 $\IndCoh$ 与 $D$-模的关系。回顾：按定义，对（局部几乎有限型的）预叠 $Y$，
\[
D\text{-}\mathrm{mod}(Y):=\IndCoh(Y_{\mathrm{dR}}).
\]
映射 $\iota\colon Y\to Y_{\mathrm{dR}}$ 给出连续且保守的函子
\[
\iota^{!}\colon D\text{-}\mathrm{mod}(Y)=\IndCoh(Y_{\mathrm{dR}})\to \IndCoh(Y),
\]
我们记之为 $\oblv_{D\text{-}\mathrm{mod}(Y)}$。
对一般预叠，$\oblv_{D\text{-}\mathrm{mod}(Y)}$ 不必有左伴随。但当 $Y$ 是代数叠时，
$\oblv_{D\text{-}\mathrm{mod}(Y)}$ 有左伴随
\[
\ind_{D\text{-}\mathrm{mod}(Y)}\colon \IndCoh(Y)\to D\text{-}\mathrm{mod}(Y)
\]
\srcpage{20}
（见 [DG13, 6.3]）。命题~\ref{prop:comp-gen-adjoint} 与命题~\ref{prop:indcoh-coh}
合起来给出

\begin{proposition}
\label{prop:qca-dmod-cg}
设 $Y$ 为 QCA 代数叠。则 $D\text{-}\mathrm{mod}(Y)$ 紧生成。
\end{proposition}

\subsection{截断性与可截断性}
\label{sec:trunc}

本节的目标是综述「$D\text{-}\mathrm{mod}(\Bun_G)$ 紧生成」的论证。若 $\Bun_G$ 是 QCA，
则由命题~\ref{prop:qca-dmod-cg} 即得结论。可惜 $\Bun_G$ 并非 QCA，因为它不拟紧。
不过它是局部 QCA 的，即其域值点的自同构群都是仿射的。

设 $Y$ 为局部 QCA 代数叠。记 $\mathrm{Open}(Y)$ 为拟紧开子叠 $U\subset Y$ 构成的偏序集。
有
\[
D\text{-}\mathrm{mod}(Y)=\lim_{U_i\in \mathrm{Open}(Y)^{\mathrm{op}}}D\text{-}\mathrm{mod}(U_i),
\]
其中对包含 $j_{i,j}\colon U_i\hookrightarrow U_j$，过渡函子是
\[
j^{!}_{i,j}\colon D\text{-}\mathrm{mod}(U_j)\to D\text{-}\mathrm{mod}(U_i).
\]
若 $j^{!}$ 都有良定义的左伴随，就落入命题~\ref{prop:comp-gen-colim} 的情形，
从而可以断言
\[
D\text{-}\mathrm{mod}(Y)=\lim_{U_i\in \mathrm{Open}(Y)^{\mathrm{op}}}D\text{-}\mathrm{mod}(U_i)
=\colim_{U_i\in \mathrm{Open}(Y)}D\text{-}\mathrm{mod}(U_i)
\]
紧生成。可惜 $j^{!}_{i,j}$ 并不总有左伴随。

\begin{definition}[{[DG15, 定义 3.1.5]}]
\label{def:truncative}
设 $Y$ 为 QCA 代数叠。设 $i\colon Z\hookrightarrow Y$ 为闭子叠，$j\colon U\hookrightarrow Y$
为其补开子叠。若 $i_*$ 有良定义的左伴随
\[
i^{*}\colon D\text{-}\mathrm{mod}(Y)\to D\text{-}\mathrm{mod}(Z),
\]
等价地，若 $j^{!}$ 有良定义的左伴随
\[
j_{!}\colon D\text{-}\mathrm{mod}(U)\to D\text{-}\mathrm{mod}(Y),
\]
则称 $Z\hookrightarrow Y$ 是\emph{截断的}（truncative），并称 $U\hookrightarrow Y$ 是
\emph{余截断的}（co-truncative）。
\end{definition}
\srcpage{21}
上述定义中两个条件（以及另外几个条件）等价的证明见 \cite[命题 3.1.2]{DG15}。

\begin{definition}[{[DG15, 定义 3.8.2]}]
\label{def:truncative-local}
现设 $Y$ 为局部 QCA 代数叠。称闭子叠 $i\colon Z\hookrightarrow Y$ 是\emph{截断的}，
若对每个拟紧开子叠 $\overset{\circ}{Y}\subset Y$，交
\[
Z\cap \overset{\circ}{Y}\hookrightarrow \overset{\circ}{Y}
\]
都是截断的。同样，称开子叠 $j\colon U\hookrightarrow Y$ 是\emph{余截断的}，
若对每个拟紧开子叠 $\overset{\circ}{Y}\subset Y$，交
\[
U\cap \overset{\circ}{Y}\hookrightarrow \overset{\circ}{Y}
\]
都是余截断的。
\end{definition}

现设 $Z\hookrightarrow Y$ 仅局部闭。则可把包含分解为
\[
Z\hookrightarrow V\hookrightarrow Y,
\]
其中 $Z\hookrightarrow V$ 是闭嵌入、$V\hookrightarrow Y$ 是开嵌入。此时称
$Z\hookrightarrow Y$ 截断，若 $Z\hookrightarrow V$ 截断。
可以形式地推出：即使 $Y$ 只是局部 QCA，$j\colon U\hookrightarrow Y$ 余截断也蕴涵
$j^{!}$ 有良定义的左伴随，仍记作 $j_{!}$。
\begin{definition}[{[DG15, 定义 4.1.1]}]
\label{def:truncatable}
称局部 QCA 代数叠 $Y$ 是\emph{可截断的}（truncatable），若它能被余截断的拟紧开子叠覆盖。
等价地，$Y$ 可截断，若 $\mathrm{CTrnk}(Y)$——即 $\mathrm{Open}(Y)$ 中由余截断拟紧开子叠
构成的子偏序集——在 $\mathrm{Open}(Y)$ 中共尾。
\end{definition}
\srcpage{22}
2.1.3 小节的讨论给出

\begin{proposition}
\label{prop:truncatable}
设 $Y$ 为可截断的局部 QCA 代数叠。则 $D\text{-}\mathrm{mod}(Y)$ 紧生成。
\end{proposition}

有人可能提出异议：对 $U\in \mathrm{CTrnk}(Y)$，我们只知道 $j_{!}$ 对 $j\colon U\hookrightarrow Y$
有定义，而 2.1.3 小节的讨论需要 $j_{i,j!}$ 对 $j_{i,j}\colon U_i\hookrightarrow U_j$ 有定义。
但注意有拉回方块
\[
\begin{tikzcd}[column sep=large,row sep=large]
U_1\times_Y U_2 \arrow[r,"j_{1,2}"] \arrow[d] \arrow[dr,phantom,"\lrcorner" description] & U_2 \arrow[d,"j_2"]\\
U_1 \arrow[r,"j_1"] & Y,
\end{tikzcd}
\]
即 $U_1\cap U_2\hookrightarrow U_1$；并且由 $D$-模的基变换，
$
j^{!}_{1,2}=j^{!}_{1}j_{2*}$，
故
\[
j_{1,2!}=j_{2}^{*}j_{1!}.
\]

\subsection{收缩子叠}
\label{sec:contractive}

本小节给出一个几何性的充分条件，保证子叠截断。按 \cite[5.1.1]{DG15} 考虑下述设定。

\begin{definition}
\label{def:globally-contractive}
把 $\mathbb{A}^1$ 视为乘法幺半群，设 $W$ 为 $\mathbb{A}^1$-概形。设
$p\colon W\to S$ 为仿射的 $\mathbb{A}^1$-等变概形映射，其中 $\mathbb{A}^1$ 在 $S$ 上作用平凡。
设 $\iota\colon S\to W$ 是 $p$ 的截面，且复合
\[
W\to S\xrightarrow{\ \iota\ }W
\]
\srcpage{23}
对应于 $0\in \mathbb{A}^1$ 的作用。记 $Y:=W/\mathbb{G}_m$，
$Z:=S/\mathbb{G}_m=S\times \mathrm{pt}/\mathbb{G}_m$。有映射
\[
i\colon Z\hookrightarrow Y
\quad\text{与}\quad
\pi\colon Y\to Z,
\]
分别由 $\iota$ 与 $p$ 得到。称形如 $i\colon Z\hookrightarrow Y$ 的闭子叠为\emph{整体收缩的}
（globally contractive）。
\end{definition}

收缩原理（见 \cite[命题 5.3.2]{DG15}）说：在上述情形下有伴随
\[
\pi_*\colon D\text{-}\mathrm{mod}(W/\mathbb{G}_m)\rightleftarrows D\text{-}\mathrm{mod}(S/\mathbb{G}_m):i_*,
\]
故可把 $i^{*}$ 等同于良定义的 $\pi_*$。换言之，整体收缩的子叠自动截断。
此概念可稍作推广：

\begin{definition}[{[DG15, 5.2.1]}]
\label{def:contractive}
称局部闭子叠 $Z'\hookrightarrow Y'$ 是\emph{收缩的}（contractive），若存在交换图
\[
\begin{tikzcd}[column sep=large,row sep=large]
Z \arrow[r] \arrow[d] & Y \arrow[d]\\
Z' \arrow[r] & Y'
\end{tikzcd}
\]
使得
\begin{enumerate}
\item[(i)] $Z\to Y$ 整体收缩；
\item[(ii)] 态射 $Z\to Z'\times_{Y'}Y$ 是开嵌入；
\item[(iii)] $Z\to Z'$ 满；
\item[(iv)] $Y\to Y'$ 光滑。
\end{enumerate}
\end{definition}
\srcpage{24}
若有如上述定义中的交换图，但上水平箭头只是收缩的而非整体收缩的，则仍可断定
$Z'\to Y'$ 收缩。

\begin{proposition}[{[DG15, 推论 5.2.3]}]
\label{prop:contractive-truncative}
若局部闭子叠收缩，则它截断。
\end{proposition}

本小节最后给出一个便于验证「某局部闭子叠收缩」的判据。

\begin{proposition}[{[DG15, 引理 5.4.3]}]
\label{prop:contractive-criterion}
设 $\pi\colon W\to S$ 为叠的仿射概形态射，且 $\mathbb{A}^1$ 经 $S$-自同构作用在 $W$ 上。
设
\begin{enumerate}
\item[(i)] $W$ 的、对应于 $0\in\mathbb{A}^1$ 的 $S$-自同构可分解为 $i\circ \pi$，
其中 $i\colon S\hookrightarrow W$ 是 $\pi$ 的某个截面；
\item[(ii)] $\mathbb{G}_m$-作用在一个点上平凡（但这不蕴涵 $\mathbb{G}_m$ 在 $W$ 上关于 $S$ 的作用平凡）。
\end{enumerate}
则 $i\colon S\hookrightarrow W$ 收缩。
\end{proposition}

\subsection{\texorpdfstring{$\Bun_G$ 的分层}{G 的分层}}
\label{sec:hn}

本小节记录 $\Bun_G$ 的 Harder–Narasimhan 分层（但不给出其定义）的若干重要事实，
以及相关的组合学。$\Bun_G$ 的 HN 分层，对 $G=\mathrm{GL}_n$ 归功于 Harder–Narasimhan，
对一般约化群 $G$ 由 Schieder 在 [Sch15] 中给出。

存在拟紧局部闭子叠
\[
\Bun^{(\lambda)}_G\subset \Bun_G,
\]
称为 Harder–Narasimhan 分层，指标集为 $G$ 的有理支配余权 $\lambda\in \Lambda^{+,Q}_G$；
定义见 \cite[定理 7.4.3]{DG15}。给定子集 $S\subset \Lambda^{+,Q}_G$，
\srcpage{25}
记
\[
\Bun^{(S)}_G:=\bigcup_{\lambda\in S}\Bun^{(\lambda)}_G .
\]
诸 $\Bun^{(\lambda)}_G$ 构成 $\Bun_G$ 的分层：它们两两不交，并为整个 $\Bun_G$。
细节见 \cite[定理 7.4.3]{DG15}。

使 $\Bun^{(\lambda)}_G$ 非空的子集 $A\subset \Lambda^{+,Q}_G$ 是离散的。特别地，给定
$\theta\in \Lambda^{+,Q}_G$，子集
\[
\{\lambda\in A\mid \lambda\le \theta\}
\]
有限，故
\[
\Bun^{(\le\theta)}_G:=\bigcup_{\lambda\le \theta}\Bun^{(\lambda)}_G
\]
拟紧。

HN 分层的拓扑在 \cite[推论 7.4.11]{DG15} 中描述。要点是：若 $\lambda\le \lambda'$，
则 $\Bun^{\lambda'}_G$ 含于 $\Bun^{\lambda}_G$ 的闭包中。结合上述 $A$ 的离散性，
可知 $\Bun^{(\le\theta)}_G$ 是 $\Bun_G$ 的开子叠。

下面把注意力转向 $\Lambda^{+,Q}_G$ 本身。先注意：对任意标准抛物子群 $P$、Levi 因子 $M$，
可考虑
\[
\Lambda^{++,Q}_{G,P}:=\{\lambda\in \Lambda^{+,Q}_G\mid
\langle\lambda,\check\alpha_i\rangle=0\ \ \forall i\in \Gamma_M,\
\langle\lambda,\check\alpha_i\rangle>0\ \ \forall i\notin \Gamma_M\}.
\]
（此处及以后，对约化群 $H$，$\Gamma_H$ 记其 Dynkin 图的节点集；我们把它与 $H$ 的单根集、
单余根集等同。）称余权 $\lambda\in \Lambda^{++,Q}_{G,P}$ 为\emph{$P$-正则的}。

\begin{definition}[{[DG15, 定义 8.2.2]}]
\label{def:padmissible}
称子集 $S\subset \Lambda^{+,Q}_G$ 是\emph{$P$-可容的}（$P$-admissible），若
\begin{enumerate}
\item[(1)] 存在 $P$-正则的 $\mu$，使得
\[
S\subset \mu+\sum_{i\in \Gamma_M}\mathbb{Q}^{\ge 0}\alpha_i;
\]
\item[(2)] 若 $\lambda_1\in S$、$\lambda_2\in \Lambda^{+,Q}_G$，且
$\lambda_1-\lambda_2\in \sum_{i\in \Gamma_M}\mathbb{Q}^{\ge 0}\alpha_i$，则 $\lambda_2\in S$；
\item[(3)] 对一切 $\lambda\in S$，只要 $i\notin \Gamma_M$，就有
$\langle\lambda,\check\alpha_i\rangle>0$。
\end{enumerate}
\srcpage{26}
进而，称 $S$ \emph{收缩 $P$-可容的}（contractively $P$-admissible），若 $S$ 是 $P$-可容的，
且对一切 $\lambda\in S$，只要 $i\notin \Gamma_M$，就有
\[
\langle\lambda,\check\alpha_i\rangle>2g-2.
\]
\end{definition}

\subsection{\texorpdfstring{$\Bun_G$ 的可截断性}{G 的可截断性}}
\label{sec:truncatable}

本小节最终断定 $D\text{-}\mathrm{mod}(\Bun_G)$ 紧生成。做法是把 $\Bun_G$ 写成拟紧余截断开子叠的并，
并说明这些开子叠的补是收缩的、因而截断的。

\begin{proposition}[{[DG15, 定理 9.1.2]}]
\label{prop:open-substack-cotrunc}
只要对一切单余根 $\check\alpha_i\in \Gamma_G$ 都有
\[
\langle\theta,\check\alpha_i\rangle\ge 2g-2,
\]
开子叠
\[
\Bun^{(\le\theta)}_G=\bigcup_{\lambda\le \theta}\Bun^{(\lambda)}_G
\]
就是余截断的。
\end{proposition}

\begin{proposition}[{[DG15, 引理 11.3.4]}]
\label{prop:square-M-P}
设 $S\subset \Lambda^{+,Q}_G$ 为收缩 $P$-可容子集。则有交换图
\[
\begin{tikzcd}[column sep=huge,row sep=large]
\Bun^{(S)}_M \arrow[r,"p^{(S)}_{P^-}"] \arrow[d,"p^{(S)}_P"'] &
\Bun^{(S)}_{P^-} \arrow[d,"\iota^{(S)}_{P^-}"]\\
\Bun^{(S)}_P \arrow[r,"\iota^{(S)}_P"'] & \Bun_G,
\end{tikzcd}
\]
\srcpage{27}
使得
\begin{enumerate}
\item[(a)] $\Bun^{(S)}_M\to \Bun^{(S)}_P\times_{\Bun_G}\Bun^{(S)}_{P^-}$ 是开嵌入；
\item[(b)] $\Bun^{(S)}_M\to \Bun^{(S)}_P$ 满；
\item[(c)] $\Bun^{(S)}_{P^-}\to \Bun_G$ 光滑；
\item[(d)] $\Bun^{(S)}_M\to \Bun^{(S)}_{P^-}$ 是仿射概形态射
$\Bun^{(S)}_{P^-}\to \Bun^{(S)}_M$ 的截面，故为闭嵌入；
\item[(e)] $\Bun^{(S)}_P\to \Bun_G$ 是到 $\Bun^{(S)}_G\subset \Bun_G$ 上的同构。
\end{enumerate}
\end{proposition}

此外（\cite[命题 11.3.7]{DG15}），存在 $\Bun_{P^-}\to \Bun_M$ 上的 $\mathbb{A}^1$-作用，
满足命题~\ref{prop:contractive-criterion} 的假设，且该作用保持 $\Bun^{(S)}_{P^-}$。
此作用的构造见 \cite[11.2]{DG15}。于是 $\Bun^{(S)}_M\to \Bun^{(S)}_{P^-}$ 收缩；
由收缩性的稳定性可得 $\Bun^{(S)}_P=\Bun^{(S)}_G\to \Bun_G$ 也收缩。

我们想说明：只要 $\theta$ 在每个方向上都「足够深」（如命题~\ref{prop:open-substack-cotrunc}
所述），$\Bun^{(\le\theta)}_G$ 就是 $\Bun_G$ 中的余截断开子叠。换言之，要验证
$\Bun^{(\le\theta)}_G\cap \overset{\circ}{Y}$ 对一切拟紧开子叠
$\overset{\circ}{Y}\subset \Bun_G$ 都在 $\overset{\circ}{Y}$ 中余截断。
其实只需对覆盖 $\Bun_G$ 的一族拟紧开子叠验证，方便的一族是
\[
\{\Bun^{(\le\theta')}_G\}_{\theta'\ge \theta}.
\]
\srcpage{28}
现注意
\[
\Bun^{(\le\theta')}_G\setminus \Bun^{(\le\theta)}_G
=\bigcup_{\lambda\not\le \theta,\ \lambda\le \theta'}\Bun^{(\lambda)}_G,
\]
并把它写成有限个 $\Bun^{(S)}_G$ 之并，其中 $S\subset \Lambda^{+,Q}_G$ 收缩 $P$-可容
（$P$ 可变动）。按 \cite[9.3.2]{DG15}，令
\[
(\theta,\theta']:=\{\lambda\in \Lambda^{+,Q}_G\mid \lambda\not\le \theta,\ \lambda\le \theta'\}.
\]
对每个 $\lambda\in(\theta,\theta']$，我们说明如何造出抛物子群 $P_\lambda$ 与含 $\lambda$ 的
收缩 $P_\lambda$-可容子集 $S_\lambda\subset(\theta,\theta']$。
令 $P_\lambda$ 为使 Levi 因子 $M_\lambda$ 对应子集
\[
\{i\in \Gamma_G\mid \langle\lambda,\check\alpha_i\rangle\le 2g-2\}
\]
的抛物子群。再取 $S_\lambda$ 为收缩 $P_\lambda$-可容子集
\[
\{\mu\in \Lambda^{+,Q}_G\mid \lambda-\mu\in \sum_{i\in \Gamma_{M_\lambda}}\mathbb{Q}_{\ge 0}\alpha_i\}.
\]
注意 $A\cap(\theta,\theta']$ 有限，故
\[
\Bun^{(\le\theta')}_G\setminus \Bun^{(\le\theta)}_G
=\bigcup_{\lambda\in A\cap(\theta,\theta']}\Bun^{(S_\lambda)}_G
\]
是有限个截断子叠之并。由于有限个截断局部闭子叠之并仍截断，命题~\ref{prop:open-substack-cotrunc}
得证。

显然诸 $\Bun^{(\le\theta)}_G$ 覆盖 $\Bun_G$，故命题~\ref{prop:open-substack-cotrunc} 给出

\begin{proposition}
\label{prop:bun-truncatable}
$\Bun_G$ 可截断。
\end{proposition}

由于 $\Bun_G$ 局部 QCA，由命题~\ref{prop:truncatable}，
$D\text{-}\mathrm{mod}(\Bun_G)$ 紧生成。
\srcpage{29}
